\section{Discussion}
\label{sec:discussion}

\subsection{Principal findings}

A selective prediction policy calibrated at one hospital did not transport. At
the external site, with clinical context fully intact, selective error rose by
\num{0.047} for probability late fusion and \num{0.038} for feature fusion,
each several times the materiality threshold and each with an interval
comfortably excluding zero. Nothing had to be manipulated to produce this: the
degradation is a property of moving between institutions.

Two of the three remaining hypotheses came back contrary to prediction, and one
of those reversals is decisive. We set out expecting multimodal models to carry
greater exposure, on the reasoning that a model which gains from context has
more to lose when context degrades. The opposite held. Both multimodal models
carried \emph{smaller} transfer gaps than the image-only control
(\cref{tab:h3}), with intervals lying entirely below zero and magnitudes above
materiality. The compound interaction was likewise not super-additive: context
degradation cost less at the external site than at the source, not more.

The transport result did not survive its pre-specified label-source
sensitivity, for reasons we take up in \cref{sec:s1-discussion}. It is
established for the primary label endpoint and is not established as robust to
label source.

We report all of this as it came out. The protocol was frozen, hash-manifested
and validated before any model existed, and the analysis code was written and
tested before these numbers were computed. The value of that discipline is
precisely that it makes a contradicted prediction reportable rather than
negotiable, and that it forced a sensitivity whose failure we would otherwise
have had no occasion to discover.

\subsection{The policy fails before the models do}

The pattern across \cref{tab:selective-risk} and the coverage results is
consistent: institutional transfer degrades selective risk substantially, while
the abstention mechanism holds its operating point --- until context is
manipulated, at which point coverage moves sharply.

Under intact context, no model breached the prespecified coverage-drift
criterion at either site, and feature fusion sat within \num{0.005} of its
\SI{80}{\percent} target at a hospital it had never seen. All four breaches
arose under intervention. The clearest is mechanical rather than subtle: with
every context replaced by one placeholder, the text-only model emits a constant
prediction, every study clears the confidence cutoff at once, and coverage goes
to \SI{100}{\percent}. A system that abstains on nothing has stopped being a
selective classifier while continuing to report a confidence.

This distinction has practical weight. A deployment that monitored coverage
alone would have seen an operating point that looked correct at the new site
while error rose by a third to two thirds. A deployment that monitored accuracy
alone would have missed the condition under which the abstention mechanism
degenerates. Neither signal is sufficient by itself.

\subsection{Why the multimodal reversal is worth taking seriously}

The reversal is not a small effect near a threshold. Both intervals exclude zero
on the side opposite the prediction, and the advantage of feature fusion over
the image-only control \emph{widens} across sites, from \num{0.034} to
\num{0.069}.

It is also not explained by the text arm being independently robust. The
text-only model had the \emph{largest} transfer gap of the four
(\num{+0.076}). Both unimodal models transported worse than both multimodal
ones. Whatever is happening, it is not that one modality is site-invariant and
carries the other.

We offer two candidate explanations and label them as conjecture, since the
design was not built to separate them. The first is that jointly trained fusion
constrains the image pathway from relying on acquisition characteristics
specific to the source institution, in the way an auxiliary signal can act as a
regulariser. The second is that the two modalities fail on partially disjoint
subsets of studies, so a fused confidence ranks errors more stably under shift
than either alone even when both are individually degraded. Distinguishing these
would require ablations the present protocol does not contain, and we do not
claim to have done so.

What can be said without conjecture is narrower and still useful: on this pair
of hospitals, with these labels and this policy, adding pre-diagnostic context
did not increase the cost of institutional transfer, and the burden of proof for
the opposite claim now sits with anyone asserting it.

\subsection{Sub-additive rather than compounding shift}

The interaction estimates were near zero or negative, so removing or corrupting
context cost less at the external site than at the source. One explanation is
available from the design itself and belongs in the interpretation rather than
being left implicit.

The informativeness rule that admits a study to the intervention counts
effective tokens. As \cref{sec:limitations} states, it detects emptiness and
brevity, not vagueness. Studies at the two sites that both qualify as
informative need not carry equal diagnostic signal, and if external context is
on average thinner in content while passing the same token threshold, then
removing it necessarily costs less. That is a limitation of the state
assignment, not evidence about compound shift, and it is the reading we favour
over any claim that compound shift is benign.

Misaligned context remained worse than absent context in three of four
instances (\cref{tab:h4}), which is the direction predicted. Wrong information
is not equivalent to no information, and a system that treats a populated
context field as evidence of a usable one will pay for that.

\subsection{What the label-source sensitivity establishes, and what it does not}
\label{sec:s1-discussion}

H1 does not survive its own pre-specified sensitivity. Under findings-derived
labels the cross-site difference reverses, and we report that rather than
qualifying it away. The claim that selective reliability fails to transport is
therefore established for the primary label endpoint and \emph{not} established
as robust to label source.

The reason is measurable and was anticipated. The frozen harmonisation plan
designated the findings endpoint a sensitivity rather than a primary on the
stated grounds that it is sparse, and required label distributions be compared
site-to-site before the cross-site experiment. That comparison
(\cref{tab:comparability}) shows the two findings label sets differ by a factor
of \num{4.9} in how often a cell carries a label at all, and by \num{35}
percentage points in what fraction of labelled cells are positive, against
\num{1.1} points under the primary endpoint. A risk computed over
\SI{9}{\percent} of cells that are \SI{69}{\percent} positive is not commensurable
with one computed over \SI{45}{\percent} of cells that are \SI{33}{\percent}
positive, and a difference between two such quantities is not a transport
estimate.

That reading is constrained rather than convenient, and two observations close
it off from the alternatives. First, the divergence separates exactly by contrast
type: every contrast comparing across sites moves, every contrast comparing
models or conditions within a site holds, in all eight instances. Had the
sensitivity merely been noisier, the disagreement would not have partitioned
along the one axis that label non-comparability predicts.

Second, and more directly, the image-only control reverses too. M1 consumes no
text, and its external cohort is identical between the two runs, so it emits
bit-identical predictions under both label sources. Its transfer gap still flips
sign. A difference computed from predictions that did not change cannot have
changed because of the model, the images, or the context; the labels are the only
thing that moved. This is not an inference we could have declined to draw.

Two consequences follow. H3 and H4, which never compare across sites within a
single estimate, carry their label-source robustness. H1 and H2 do not, and the
strength of those claims should be read accordingly.

The broader point is methodological, and we think it is the more useful one.
Cross-site evaluation presupposes commensurable labels, and label sets bearing
the same name at two institutions need not be commensurable. Ours were not, and
the divergence is large enough to reverse a sign. Any study comparing risk
across institutions inherits this requirement whether or not it checks it; we
checked it only because the protocol demanded the check in advance. This is a
concrete instance of the argument that cross-domain failure in chest radiography
is driven by label shift~\cite{cohen2020limits}, arriving here not as a
competing explanation for model behaviour but as a constraint on what can be
compared at all.

\subsection{Relation to prior work}

Our source-site results are consistent with the finding that clinical history
improves chest-radiograph classification~\cite{yang2025mcxnet}: feature fusion
attained the lowest selective error under every condition at both sites. Where
we depart is in what that improvement implies. That work observes naturally
context-free images within one institution; we intervene on context under a
frozen rule and evaluate across institutions, and find that the benefit does not
convert into transfer fragility.

The dominance of institutional over contextual degradation is consonant with
the argument that cross-domain failure in chest radiography is driven by label
shift rather than image shift~\cite{cohen2020limits}, and with the observation
that calibration and cross-center generalisation remain open after multi-center
benchmarking~\cite{cxrlt2026}. Our contribution to that line is to show the
failure persists when the quantity being transferred is an explicit abstention
policy rather than a classifier alone, and to quantify it at a fixed operating
burden.

Our reversal also has a precedent worth naming. Comparing architectures for
robustness to acquisition view type, model rankings reversed entirely on a
second dataset, with the model best on internal validation becoming the worst
externally~\cite{farquhar2026acquisition}. Our H3 result is a reversal of the
same kind and in the opposite direction: the ordering we predicted from
source-site behaviour did not survive the move, and the models we expected to be
most exposed proved least. Two findings of this shape, on different shift axes
and different architectures, suggest a general caution rather than a
peculiarity of either study. Internal ranking is not evidence about external
ranking, and reporting one as though it licensed the other is a recurring error
this literature has yet to stop making.

The missing-modality literature treats absent inputs as an obstacle to engineer
around~\cite{wang2025missing,wang2025moehealth}. Our results suggest a
reordering of priorities for deployment: on this evidence, the modality
robustness problem was less costly than the institutional transfer problem, and
architectural effort spent on the former would not have addressed the latter.

\subsection{Methodological observations}

Two findings arose from constructing the protocol and generalise past it.

A substantial minority of source reports embed a preliminary radiologist
interpretation inside the report body. Any pipeline that takes text preceding
the findings section as pre-diagnostic context ingests post-diagnostic text for
those studies. This is a silent contamination: it inflates apparent context
value without producing any visible failure, and it affects any study using this
corpus in this way, not only ours.

Separately, the decision rule this literature commonly adopts --- requiring both
that an interval exclude zero and that the point estimate reach a materiality
threshold --- caps power at \SI{50}{\percent} when the true effect equals that
threshold, for any sample size. Studies adopting such a rule should power
against a minimum detectable effect and state materiality as the smallest effect
worth acting on, not as the effect they are powered to detect.

\subsection{Implications for deployment}

A frozen abstention threshold should not be assumed to transport. On this
evidence it does not, and the failure is invisible to coverage monitoring under
normal operation, because coverage was well behaved at the external site
precisely where error was not.

Re-calibrating at the target site would have concealed this. A study that
re-fitted thresholds externally would have reported an operating point restored
and a policy that transports, when what had actually been demonstrated is that
the method can be re-tuned. Whether the shipped artifact survives the move is a
different question, and answering it requires freezing the artifact.

\subsection{Limitations bearing on interpretation}

\cref{sec:limitations} states the full set; three bear directly on how these
conclusions should be read.

Single-seed training bears on the H3 reversal specifically. Its magnitude,
\numrange{0.027}{0.035}, is not far above what training variation alone might
produce, and the bootstrap does not speak to that source of uncertainty. Two
features of the result argue against a seed artefact: the direction holds across
two independently constructed fusion models, and the gap widens across sites
rather than staying flat. Neither is decisive, and a replication with repeated
seeds would settle it.

Note also that H1, the confirmed primary, is insulated from this concern in a
way H3 is not. H1 compares the same weights evaluated at two sites, so whatever
idiosyncrasies a particular training run produced are carried to both and
largely cancel in the difference. H3 compares two different models, where no
such cancellation occurs.

The two sites differ in more than institution. Their images have different
compression histories at origin, their permitted context arrives through
different mechanisms, and their label prevalences differ. The design isolates
context by intervention within each site, but the institutional contrast remains
a comparison of two hospitals rather than of institution alone.

Finally, this is one pair of hospitals. Nothing here establishes how large a
transfer gap to expect elsewhere, only that a frozen policy failed to transport
between these two by a margin that would matter clinically.

\subsection{Limitations}
\label{sec:limitations}


Several limitations are known independently of the confirmatory result and are
stated here rather than deferred.

\paragraph{Effective token count measures brevity, not vagueness.}
The informativeness rule that separates informative from low-information context
counts tokens carrying alphanumeric content after de-identification placeholders
are removed. It therefore detects emptiness and brevity. A genuinely vague but
lengthy indication --- ``evaluate'' followed by boilerplate --- scores as
informative. This is a measurement limitation of the state assignment, not of
the controlled interventions, which do not depend on it beyond eligibility.

\paragraph{The primary evaluation partition is not an official split.}
The official test split of the source dataset yields too few patients within the
study cohort to support the primary comparison, so a patient-disjoint partition
was carved from the training split using patient identifiers and a frozen seed
alone. Primary estimates are consequently not directly comparable to published
official-split baselines. The official validation and test splits are retained
and reported as secondary confirmation.

\paragraph{The transport result is not robust to label source.}
H1 is established under the primary label endpoint and reverses under the
pre-specified sensitivity endpoint. The two findings-derived label sets proved
non-commensurable across sites, differing by a factor of five in labelling
density, so the sensitivity cannot arbitrate the cross-site contrast in either
direction. What this leaves is a transport failure demonstrated for one label
endpoint, with its label-source robustness untested rather than refuted.
Establishing it would require a label endpoint that is dense and comparable at
both institutions, which neither available endpoint is.

\paragraph{Two training seeds, and only for two of the four models.}
The bootstrap quantifies uncertainty in \emph{evaluation}, not variability
across training runs. A second seed was run for the image-only and
feature-fusion models and reproduced every verdict
(\cref{sec:seed-replicate}), but two runs bound seed variability loosely: point
estimates moved by up to a third of their own magnitude while keeping their
sign and their verdict. Differences between models smaller than the materiality
threshold remain unestablished against seed-to-seed variation. The text-only
model was not retrained, so neither it nor the late-fusion model that consumes
it has been replicated.

\paragraph{Source images are lossy, external images are lossless.}
The source dataset publishes JPEG images and the external dataset publishes PNG.
Both were reduced to the model input resolution by a single identical bilinear
step, so no additional resampling asymmetry was introduced, but the compression
history of the two sites differs at origin. This difference tracks the site
variable and cannot be removed by preprocessing choices.

\paragraph{One external image is unusable at source.}
One frontal image in the external dataset matches the publisher's recorded
checksum yet fails to decode: its PNG header parses while an IDAT chunk fails
its cyclic redundancy check. The damage is upstream and no retransfer repairs
it. That study is excluded, and the exclusion is reported rather than absorbed.

\paragraph{Studies whose only images lack a view label are excluded.}
Images carrying no view label are treated as non-frontal, which is the
conservative reading, since assuming them frontal would admit views the protocol
excludes. Such exclusions may not be missing at random.

\paragraph{Section boundaries derive from a curated header vocabulary.}
Reports whose headers fall outside that vocabulary are excluded rather than
parsed heuristically, which trades coverage for the guarantee that no
post-diagnostic text enters a pre-diagnostic field.

\paragraph{No sociodemographic subgroup or fairness analysis was performed.}
This study reports no patient demographics and no subgroup performance
estimates, and applied no fairness-specific method. The omission was not a
considered exclusion: the protocol specified cohorts, endpoints, contrasts and
materiality without naming sociodemographic strata, so no such analysis was
frozen in advance, and adding one after the external results were inspected
would be post-hoc exploratory by the protocol's own rule.

The gap matters here more than it would in a study reporting only aggregate
accuracy, because it interacts with the central finding. If a frozen abstention
policy transports poorly between institutions, there is no reason to assume it
transports uniformly across patient groups within one. A policy that deferred
disproportionately on one group would be a fairness failure invisible to every
metric reported here --- the same silence that makes coverage monitoring
inadequate for institutional transfer would hide it. Both datasets publish
demographic fields sufficient to test this, and doing so is the natural
companion to the third site identified in \cref{sec:future-work}.

\subsection{Future work}
\label{sec:future-work}

The most informative next step is not a better model but a third site. Two sites
give a difference; three begin to give a distribution, and would show whether
the multimodal transfer advantage observed here is a property of these
institutions or of multimodal fusion.

Separately, the sub-additivity result motivates an informativeness rule that
measures content rather than length. A rule that admitted studies on the
strength of what their context says, rather than how much of it there is, would
make the interaction estimate interpretable in a way the present one is not.
